\section{Introduction}

Tracks is a logic puzzle defined on a rectangular grid. Two terminals are
prescribed on the board, and the objective is to construct a railway route
connecting them. The solution must satisfy two families of conditions
simultaneously. On the one hand, the route must be locally and globally
consistent: neighboring track pieces must match, the structure must not branch,
and the final railway must form a single valid route between the prescribed
terminals. On the other hand, the number of cells occupied by the route in each
row and in each column must match the clues written on the boundary of the
grid.

From the perspective of operations research, Tracks is naturally interpreted
as a discrete feasibility problem on a graph. The combinatorial difficulty does
not only come from counting constraints, but also from the global structure of
the route. Local consistency constraints are not sufficient, since they allow
disconnected components or isolated cycles. A rigorous model must therefore
combine binary selection variables, degree conditions, and explicit
connectivity constraints. This places the puzzle in the same mathematical
family as several classical discrete-optimization problems in which a subset
of vertices must satisfy both local incidence constraints and a global
connectedness property.

This structure makes Tracks a compact case study for graph modeling. Each cell
of the puzzle becomes a vertex, each possible continuation between
orthogonally adjacent cells becomes an edge, and the final railway must be a
simple path. Graph-theoretic notions such as degree, path, cycle, connected
component, cut, and flow therefore provide a precise language for expressing
the rules of the game. The use of network-flow ideas is particularly natural:
flow can certify that every selected cell is reachable from a terminal, which
is exactly the property that local track-piece compatibility cannot guarantee
on its own \cite{ahuja1993,validi2022,miyazawa2021}.

The project follows this mathematical viewpoint from the model to the
implementation. The grid is first represented as a graph, then the puzzle
rules are translated into a mixed-integer linear formulation, and finally the
formulation is implemented in Python using a standard MILP backend. The same
objects also support parsing, validation, generation, dataset solving, and
visualization. The purpose of the report is therefore not only to present a
solver, but also to show how a visually simple puzzle can be handled as a
complete graph-based feasibility problem.
